\chapter*{Prakata}
\addcontentsline{toc}{chapter}{Prakata}

Pada semester musim dingin 2024/25 saya mulai belajar kecerdasan buatan di Linz. Semester
pertama langsung padat: computer vision, jaringan saraf, dan LSTM yang diajarkan oleh
penemunya sendiri. Setiap kuliah meninggalkan setumpuk slide. Dua tahun kemudian, map di
laptop berisi ratusan PDF.

Slide bagus untuk mengikuti kuliah, tetapi kurang enak dibaca ulang. Kalimatnya terpotong,
rumus muncul tanpa cerita, dan hubungan antar mata kuliah jarang terlihat. Padahal hubungan
itulah yang paling menarik. Konvolusi di kuliah computer vision ternyata saudara dekat
stensil beda hingga di kuliah mekanika fluida. Model difusi di kuliah machine learning lanjut
ternyata bersandar pada persamaan Fokker--Planck dari fisika statistik. Kalman filter di kuliah
model probabilistik ternyata cara yang sama yang dipakai ahli cuaca untuk menggabungkan
simulasi dengan pengukuran.

Buku ini adalah catatan perjalanan itu. Saya menulis ulang isi kuliah dengan bahasa sendiri,
menambahkan gambaran yang dulu membantu saya memahaminya, lalu merujuk literatur yang lebih
luas dari sekadar slide. Harapannya, siapa pun yang membaca, entah mahasiswa, insinyur, atau
orang yang sekadar penasaran, bisa mengikuti ceritanya dari satu neuron sampai model yang
meniru aliran fluida.

Di tengah perjalanan, satu pertanyaan semakin sering muncul di kepala saya: bagaimana machine
learning bisa membantu simulasi fisika, dan bagaimana fisika bisa membuat machine learning lebih
bisa dipercaya? Pertanyaan itu menjadi benang merah buku ini. Hampir setiap bab ditutup dengan
melihat materinya dari sudut simulasi, dan Bagian III membahas dunia \emph{physics-informed
machine learning} secara utuh: PINN, neural operator, simulator berbasis graf, penemuan
persamaan dari data, dan model hibrida.

Gaya penulisannya sengaja dibuat seperti catatan yang dibaca, bukan buku latihan. Setiap
gagasan besar diberi satu gambaran dari kehidupan sehari-hari sebelum rumusnya muncul. Analogi
tidak pernah sempurna, dan kalau sebuah analogi mulai menyesatkan, teks akan mengatakannya.
Contoh hitungan kecil ditulis di dalam paragraf, supaya pembaca bisa mengikutinya dengan
pensil kalau mau. Istilah teknis dibiarkan dalam bahasa Inggris dan dicetak miring saat pertama
muncul, karena istilah itulah yang dipakai di kuliah dan di makalah.

Setiap bab diakhiri catatan singkat tentang mata kuliah asal dan bacaan utamanya. Slide dan
soal asli tetap milik para pengajarnya dan tidak disertakan. Yang ada di sini adalah pemahaman
saya atas materi itu, jadi kekeliruan yang tersisa adalah kekeliruan saya.

\bigskip
\hfill Linz, Oktober 2026
